\section*{Chapter \ref{ch:m1:zero-day-options-and-retail-flow} --- Zero-Day Options, Weeklies and Retail Options Flow}

\begin{solution}{exo:m1:zero-day-options-and-retail-flow:1}
$0.8 \times 16\%/\sqrt{252} = 0.81\%$ of the index; at 32\%, 1.61\%. (The exact
value at 16\% is 0.80\%.)
\end{solution}

\begin{solution}{exo:m1:zero-day-options-and-retail-flow:2}
Hedge: short $0.50 \times 500 \times 100 = 25\,000$ shares. Linear estimate for
$+1\%$ (one dollar): \omterm{def:m1:zero-day-options-and-retail-flow:gamma}{delta} rises by 0.396, that is $0.396 \times 50\,000 =
19\,800$ shares (19\,790 unrounded). The desk is long \omterm{def:m1:zero-day-options-and-retail-flow:gamma}{gamma}: it \emph{sells}
them as the share rises.
\end{solution}

\begin{solution}{exo:m1:zero-day-options-and-retail-flow:3}
$\sqrt{21} = 4.6$; $\sqrt{390/15} = 5.1$. \omterm{def:m1:zero-day-options-and-retail-flow:gamma}{Gamma} multiplies by about twenty-three
between a month out and the last quarter of an hour.
\end{solution}

\begin{solution}{exo:m1:zero-day-options-and-retail-flow:4}
\omterm{def:m1:zero-day-options-and-retail-flow:gamma}{Delta} goes from 0.50 to 0.84: sell $0.34 \times 50\,000 = 16\,900$ shares
(16\,873 exactly), against 19\,790 by the linear estimate. \omterm{def:m1:zero-day-options-and-retail-flow:gamma}{Gamma} is the slope
of \omterm{def:m1:zero-day-options-and-retail-flow:gamma}{delta} at 100; at 101, one dollar being 1.0 standard deviations of the
remaining day, the option is already well in the money and its \omterm{def:m1:zero-day-options-and-retail-flow:gamma}{gamma} has
fallen: \omterm{def:m1:zero-day-options-and-retail-flow:gamma}{delta} saturates at one.
\end{solution}

\begin{solution}{exo:m1:zero-day-options-and-retail-flow:5}
$2.3 \times 10^6 \times 100 \times 6\,000 = \$1.38$ trillion of notional. Hedging
flow is \omterm{def:m1:zero-day-options-and-retail-flow:gamma}{delta} and \omterm{def:m1:zero-day-options-and-retail-flow:gamma}{gamma}, not notional: most of these options are out of the
money with small \omterm{def:m1:zero-day-options-and-retail-flow:gamma}{deltas}, and a hedger trades only the \emph{change} in \omterm{def:m1:zero-day-options-and-retail-flow:gamma}{delta}.
And it is the dealers' \emph{net} position that is hedged: customers both buy
and sell the same strikes, contracts are opened and closed many times in the
day, so the net is a small fraction of the volume. Two or three orders of
magnitude separate the two numbers.
\end{solution}

\begin{solution}{exo:m1:zero-day-options-and-retail-flow:6}
Short \omterm{def:m1:zero-day-options-and-retail-flow:gamma}{gamma}: it buys shares as the price rises 3\% in the last hour (its short
calls' \omterm{def:m1:zero-day-options-and-retail-flow:gamma}{delta} is rising) and sells them as the price falls back the next
morning. It bought high and sold low: a hedging loss roughly proportional to
\omterm{def:m1:zero-day-options-and-retail-flow:gamma}{gamma} times the square of the move. In exchange it collected the options'
premium, that is, \omterm{def:m1:volatility-first-contact:iv}{implied volatility}; it earns if the realised moves it has
to chase are smaller than those the premium paid for.
\end{solution}

\begin{solution}{exo:m1:zero-day-options-and-retail-flow:7}
0.77, 1.00 and 1.42. For $r_t = \varepsilon_t + f\,r_{t-1}$ the sum of returns
over a long horizon is approximately $\sum\varepsilon_t/(1-f)$, so the ratio is
$1/(1-f)$: 0.77 for $f = -0.3$ and 1.43 for $f = +0.3$. The asymmetry is that
of $1/(1-f)$: the ratio diverges as $f \to 1$.
\end{solution}

\begin{solution}{exo:m1:zero-day-options-and-retail-flow:8}
(i) Same-day options opened and closed during the session are not in last
night's \omterm{def:m1:futures-contracts-and-exchanges:oi}{open interest} at all; the number describes longer-dated positions.
(ii) The sign is assumed: if customers sell calls (overwriting) and also sell
puts (premium collection), dealers are long both and the true figure is
larger; if customers buy calls, as retail does, dealers are short them and
the sign flips. (iii) \omterm{def:m1:zero-day-options-and-retail-flow:gamma}{Gamma} near expiry depends violently on where the index
is relative to each strike and on the hour: a number computed at 09:00 for
``today'' is stale by 10:00. One may add that customer flow in these options
is often two-sided, so the net dealer position is far smaller than \omterm{def:m1:futures-contracts-and-exchanges:oi}{open
interest} or volume suggest.
\end{solution}

\begin{solution}{pb:m1:zero-day-options-and-retail-flow:1}
\textbf{1.} 2 million shares, \$1.2 billion.
\textbf{2.} Short 20\,000 calls with a \omterm{def:m1:zero-day-options-and-retail-flow:gamma}{delta} of 0.50: long 1 million shares.
\textbf{3.} $16\% \times \sqrt{60/98\,280} = 0.40\%$.
\textbf{4.} $0.168 \times 6 \times 2\,000\,000 = 2.02$ million shares per 1\%.
Short \omterm{def:m1:zero-day-options-and-retail-flow:gamma}{gamma}: the dealer \emph{buys} as the price rises and sells as it falls.
\textbf{5.} It exceeds the whole notional, which is impossible: \omterm{def:m1:zero-day-options-and-retail-flow:gamma}{delta} can only
go from 0.50 to 1. The linear number is valid for moves small compared with
one standard deviation, 0.40\%, and a 1\% move is two and a half of them.
\textbf{6.} $0.50 + 472\,000/2\,000\,000 = 0.74$ and $0.50 + 987\,000/2\,000\,000 =
0.99$.
\textbf{7.} After $+0.25\%$, six tenths of a standard deviation, \omterm{def:m1:zero-day-options-and-retail-flow:gamma}{delta} has
already covered half of its remaining range; beyond one standard deviation
there is little left to hedge.
\textbf{8.} $2.02/0.792 = 2.55 = \sqrt{390/60}$; $4.04/2.02 = 2.0 =
\sqrt{60/15}$.
\textbf{9.} An average hour trades $60/6.5 = 9.2$ million shares: the linear
number is 22\% of it, for one dealer and one strike.
\textbf{10.} At 600.30 the calls are just in the money and the dealer is long
most of the 2 million shares; at 599.70 it should hold few. With five
minutes left each crossing of the strike asks it to sell, then buy back, a
large part of 2 million shares.
\textbf{11.} It sells near 599.70 and buys near 600.30: up to 60 cents on the
shares it turns over, per oscillation, plus spread and impact. This is the
\omterm{def:m1:volatility-first-contact:realised}{realised volatility} it is short, paid in cash.
\textbf{12.} The premium of the 20\,000 calls, from the customers who bought
them that morning: about 0.4\% of \$1.2 billion, some \$5 million, if sold at
the open.
\textbf{13.} Assignment: if the fund closes at or near 600 it does not know
how many calls will be exercised, hence how many of its hedge shares it must
deliver (\cref{pb:m1:option-contracts-and-exchanges:1}).
\textbf{14.} It would sell above 600 and buy below: its hedging supplies
liquidity on both sides of the strike and tends to hold the price there, the
pinning of \cref{pb:m1:option-contracts-and-exchanges:1}, question 16.
\textbf{15.} They could be if dealers were heavily short \omterm{def:m1:zero-day-options-and-retail-flow:gamma}{gamma} near the
money: their buying of rises and selling of falls, growing as expiry
approaches, would feed on itself (\cref{fig:m1:zero-day-options-and-retail-flow:feedback}).
What limits it is that customers both buy and sell these options, so that
dealers' net \omterm{def:m1:zero-day-options-and-retail-flow:gamma}{gamma} is a small fraction of the volume, and of either sign.
\textbf{16.} It is measured once a day from the previous night's positions;
same-day trading opens and closes inside the session.
\textbf{17.} Many small gains and rare losses several times the premium: a
distribution with a long left tail. Position size must come from the loss on
a day when the index moves three or four daily standard deviations, not from
the average.
\textbf{18.} Volume: every expiry is a new product that needs no new
listing, attracts event-driven and retail demand, and is exclusive to the
exchange that holds the index licence.
\textbf{19.} \textbf{2.02 million shares per 1\%, that is \$1.2 billion of
shares, the whole notional.}
\textbf{20.} Its \omterm{def:m1:zero-day-options-and-retail-flow:gamma}{delta} is no longer a slowly moving number but a switch that
the underlying can flip.
\end{solution}

\begin{solution}{iq:m1:zero-day-options-and-retail-flow:1}
\omterm{def:m1:zero-day-options-and-retail-flow:gamma}{Gamma} is the rate of change of \omterm{def:m1:zero-day-options-and-retail-flow:gamma}{delta} with the underlying. Long \omterm{def:m1:zero-day-options-and-retail-flow:gamma}{gamma} (long
options): \omterm{def:m1:zero-day-options-and-retail-flow:gamma}{delta} rises with the price, so the hedger sells into rises and
buys into falls, earning from movement and paying time decay. Short \omterm{def:m1:zero-day-options-and-retail-flow:gamma}{gamma}:
the reverse; the hedger chases the market and is paid premium for it.
\iqlookfor{the direction of hedging trades and the decay-for-movement
exchange.}
\end{solution}

\begin{solution}{iq:m1:zero-day-options-and-retail-flow:2}
At the money it grows like $1/\sqrt T$ and becomes unbounded at expiry: \omterm{def:m1:zero-day-options-and-retail-flow:gamma}{delta}
jumps between zero and one at the strike. Out of the money it first rises,
as the distribution narrows towards the strike's neighbourhood, then
collapses to zero once the strike is several $\sigma\sqrt T$ away. \omterm{def:m1:zero-day-options-and-retail-flow:gamma}{Gamma}
concentrates in a band of width $\sigma\sqrt T$ around the strike.
\iqlookfor{both behaviours and the band.}
\end{solution}

\begin{solution}{iq:m1:zero-day-options-and-retail-flow:3}
$0.8\,\sigma\sqrt T$: at 16\% volatility a day is $16/\sqrt{252} \approx 1\%$,
so the straddle costs about 0.8\% of the index, 48 points at 6\,000. It breaks
even if the index closes more than 0.8\% away, about 0.8 daily standard
deviations, which has a probability of roughly 42\%.
\iqlookfor{the rule of sixteen and a probability.}
\end{solution}

\begin{solution}{iq:m1:zero-day-options-and-retail-flow:4}
Classify each option trade as customer buy or sell from the tape (capacity
codes where available, otherwise quote rule and size), accumulate
dealer-side positions by strike through the day, revalue \omterm{def:m1:zero-day-options-and-retail-flow:gamma}{gamma} at the
current spot and time. Test: regress subsequent \omterm{def:m1:volatility-first-contact:realised}{realised volatility}, or the
autocorrelation of short-horizon returns, on the estimate, out of sample,
controlling for \omterm{def:m1:volatility-first-contact:iv}{implied volatility} and time of day; check that the sign
flips when the estimate does. Be ready to find that the effect is small.
\iqlookfor{signed trades rather than \omterm{def:m1:futures-contracts-and-exchanges:oi}{open interest}; a falsifiable test.}
\end{solution}

\begin{solution}{iq:m1:zero-day-options-and-retail-flow:5}
Buy the options back and pay the last of the premium: the clean exit. Buy
other options at nearby strikes to cap the \omterm{def:m1:zero-day-options-and-retail-flow:gamma}{gamma}. Keep hedging, accepting the
cost of each crossing, with a wider hedging band to avoid being whipsawed.
Or stop hedging and carry the binary outcome, sized so that it is
affordable. The worst choice is hedging continuously with a tight band.
\iqlookfor{closing as the default; bands; sizing.}
\end{solution}

\begin{solution}{iq:m1:zero-day-options-and-retail-flow:6}
Mechanism: dealers who are net short same-day \omterm{def:m1:zero-day-options-and-retail-flow:gamma}{gamma} amplify moves, net long
damp them; \omterm{def:m1:zero-day-options-and-retail-flow:gamma}{gamma} per contract is enormous near the strike. Against a large
effect: customer flow is two-sided and much of it is spreads, so dealer net
\omterm{def:m1:zero-day-options-and-retail-flow:gamma}{gamma} is small; positions are flat by the close; hedging is done in futures
with deep liquidity. Evidence needed: dealer net \omterm{def:m1:zero-day-options-and-retail-flow:gamma}{gamma} by minute from signed
trades; its relation to intraday return autocorrelation and volatility;
event studies on days with large one-sided flow; comparison of intraday
volatility before and after daily expiries were introduced, controlling for
the level of volatility.
\iqlookfor{net versus gross; the right data.}
\end{solution}
